\chapter{Independent Classification Axes}\label{ch07:classification}

\section{Why two axes, and why they are independent}\label{ch07:sec:why}

Every classification decision in CAS-Ohm 1.0 descends from a single retained
source conflict: the Quantification Guide's ASOM mnemonic
(Starved/Aligned/Oversaturated/Misaligned) overlays semantic judgments onto
dose-response regions, while the operationalization addendum states plainly
that dose-response signatures and effect-axis outcomes are distinct and that a
flat response has multiple causes. The candidate resolves this not by choosing
one source over the other but by refusing the overlay as a matter of type
discipline: the implementation classifies two independent axes and derives any
mnemonic label afterward, as a display-only partial function that is permitted
to return nothing.

The first axis answers a question about the \emph{shape} of the local
dose-response relation at a registered test point: as the declared alpha dose
is varied over adjacent points of a sweep, does the observed flow magnitude
respond linearly at the expected slope, not at all, only up to a ceiling, or
in some other way? This is the question dose-response analysis has always
asked, and the candidate deliberately inherits the vocabulary of that
tradition rather than the vocabulary of circuit diagnosis: a region label is a
statement about a fitted local curve under a declared design, in the sense of
nonlinear dose-response regression \cite{ritz}, not a statement about what the
underlying system \emph{is} or \emph{wants}.

The second axis answers a question about the \emph{direction} of whatever
response exists: did the activation move along the registered direction
vector, against it, or somewhere else in activation space? This question is
answerable, in principle, even when the magnitude question is degenerate, and
unanswerable in practice in many situations where the magnitude question is
crisp. The two questions have different evidence requirements, different
failure modes, and different silence conditions. Collapsing them into one
label is precisely the ASOM overload the source audit preserved as a conflict,
and the whole architecture of this chapter is the refusal to perform that
collapse.

Throughout this chapter the enumerations are reproduced exactly as
implemented in the candidate's typed ontology
(\texttt{src/cas\_ohm\_v1/model.py}). The dose-response region enumeration has
seven members; the semantic outcome enumeration has four. Nothing in this
chapter asserts that any classification has been produced against a live
model: the artifact status of the entire candidate remains
\textsc{open / ring 4 / not licensed / held}, and every numeric example below
is an illustrative construction over synthetic or hypothetical sweep data, not
an empirical result.

\begin{boundary}[Ring 4 status of every classification]\label{ch07:bound:ring4}
A region label or outcome label emitted by the candidate is a statement about
recorded sweep observations under one calibration profile, produced by tested
software. It is not a scientific conclusion, a safety finding, a measurement
of attention, or a causal diagnosis. No classification in this chapter is
licensed for any representation beyond structural research use.
\end{boundary}

\section{The dose-response region axis as implemented}\label{ch07:sec:region}

\begin{definition}[Dose-response region enumeration]\label{ch07:def:region}
The implemented enumeration \texttt{DoseResponseRegion} has exactly seven
members, in declaration order: \texttt{UNASSESSED}, \texttt{CUTOFF},
\texttt{ACTIVE}, \texttt{SATURATION}, \texttt{BREAKDOWN},
\texttt{NON\_OHMIC}, \texttt{UNSTABLE}. No other region label exists in the
type system, and the classification container
(\texttt{OperatingClassification}) rejects any value that is not an exact
member of this enumeration.
\end{definition}

A structural fact about the implementation deserves emphasis before the
region-by-region treatment, because it governs how several members must be
read. The candidate distinguishes two classification modes. In
\texttt{NUMERIC\_DERIVED} mode, a factory function derives the region from at
least three sweep observations bracketing an interior index, using only the
numeric neighbors, the profile's expected resistance, and the profile's
tolerance thresholds; the factory then binds the exact observation tuple, its
canonical digest, and the classification into one assessment object that
replays deterministically. In \texttt{EXTERNAL\_UNVERIFIED} mode, a region
label supplied from outside is carried as an explicitly unverified annotation
and is forbidden from embedding any derived sweep claim. The numeric deriver,
as implemented, emits only four of the seven members: \texttt{CUTOFF},
\texttt{SATURATION}, \texttt{NON\_OHMIC}, and \texttt{ACTIVE}. The members
\texttt{UNASSESSED}, \texttt{BREAKDOWN}, and \texttt{UNSTABLE} are reachable
only as typed carriers of external or absent judgment. This asymmetry is not
an oversight; it is the candidate's honesty about what its three-point local
classifier can and cannot distinguish, and it is recorded here because a
reader of the enumeration alone would not see it.

\begin{invariant}[Numeric emission subset]\label{ch07:inv:numeric-subset}
In \texttt{NUMERIC\_DERIVED} mode the implemented classifier emits a region
in $\{\texttt{CUTOFF}, \texttt{SATURATION}, \texttt{NON\_OHMIC},
\texttt{ACTIVE}\}$ and never emits \texttt{UNASSESSED}, \texttt{BREAKDOWN},
or \texttt{UNSTABLE}. The latter three members are admissible only in
\texttt{EXTERNAL\_UNVERIFIED} mode, where the assessment may not carry
observations, an index, or an expected resistance.
\end{invariant}

The numeric derivation itself, for the record, proceeds as follows. Given
strictly increasing drives at the predecessor, current, and successor
observations, the classifier computes the backward slope, the forward slope,
the expected slope $1/r_{\mathrm{eff}}$ from the supplied expected
resistance, a slope floor equal to the expected slope times the profile's
relative tolerance, and a relative slope error equal to the larger of the two
slope deviations divided by the expected slope. Classification is then a
short guarded cascade: flow at or below the absolute tolerance yields
\texttt{CUTOFF}; a flat forward slope with a live backward slope yields
\texttt{SATURATION}; a relative slope error above the relative tolerance
yields \texttt{NON\_OHMIC}; otherwise \texttt{ACTIVE}. Every threshold in
this cascade is a profile-bound tolerance, not a universal constant, and the
same sweep classified under a different tolerance profile may land in a
different region. Region labels are therefore profile-relative twice over:
through the calibration identity that produced the observations and through
the tolerance identity that partitioned them.

\subsection{\texttt{UNASSESSED}}\label{ch07:sec:region-unassessed}

\emph{Signature.} There is none. \texttt{UNASSESSED} is the typed absence of
a region judgment: no sweep was classified, or an external annotation
declined to commit. It is the region axis analogue of a missing value that
has been made unignorable by being given a name.

\emph{Causes.} A profile may have been constructed and reviewed without any
sweep yet run; an external record may carry semantic information but no
dose-response evidence; a pipeline may deliberately defer region
classification until controls are in place. In all such cases the honest
label is \texttt{UNASSESSED}, and the type system carries it forward rather
than defaulting to anything.

\emph{What it does not imply.} \texttt{UNASSESSED} does not imply
\texttt{ACTIVE}-by-default, does not imply that the response is small, and
does not imply that assessment failed. It implies only that no admissible
classification event has occurred. Critically, an \texttt{UNASSESSED} region
participates in the no-diagnosis conditions of
Section~\ref{ch07:sec:nodiag}: since residual diagnosis requires the
\texttt{ACTIVE} region and \texttt{UNASSESSED} is not \texttt{ACTIVE},
nothing downstream fires.

\subsection{\texttt{CUTOFF}}\label{ch07:sec:region-cutoff}

\emph{Signature.} The observed flow magnitude at the classified point is at
or below the profile's absolute tolerance (\texttt{delta\_abs}): within
measurement resolution, nothing is flowing. In the implemented cascade this
test fires first, so a point that is simultaneously flat and tiny is
\texttt{CUTOFF}, not \texttt{SATURATION}.

\emph{Causes.} The candidate deliberately does not enumerate causes inside
the label, because the causes are many and the label identifies none of them.
A dose below an activation threshold; a direction vector that fails to
intersect the mechanism it was intended to probe; an intervention applied at
a layer or token position where the target computation has not yet formed or
has already been consumed; a hard upstream gate; an instrumentation fault
that zeroes the projection; a full-scale or units error that shrinks the
recorded value below tolerance. All produce the same signature.

\emph{What it does not imply.} \texttt{CUTOFF} does not imply that the
system is \emph{starved} of the driven content, that larger doses would
produce response, that the mechanism is absent, or that the intervention is
safe or ineffective in any behavioral sense. It is a statement that this
sweep, at this point, under this profile, recorded no resolvable flow. The
mnemonic reading (\textsc{starved}) is available only through the overlay of
Section~\ref{ch07:sec:asom} and only when the semantic axis independently
reads \texttt{INTENDED}, which is itself a subtle condition: a semantic
judgment about a response that is below resolution inherits all the fragility
of dividing by small numbers, which is why downstream ratio diagnostics are
censored rather than computed (Chapter~\ref{ch10:residuals} of this
monograph).

\subsection{\texttt{ACTIVE}}\label{ch07:sec:region-active}

\emph{Signature.} Flow is resolvable, both local slopes exist, and the
relative slope error against the profile's expected resistance is within the
relative tolerance. Informally: the local dose-response relation looks like
the calibrated line, at the calibrated slope, within declared error.

\emph{Causes.} An \texttt{ACTIVE} label is evidence-compatible with the
candidate's central hypothesis --- a local linear response region --- but it
is equally compatible with a locally linear stretch of a globally nonlinear
curve, with a coincidental slope match at one interior point, or with
compensating errors in the full-scale and direction conventions. Three
points cannot distinguish these; the label does not pretend to.

\emph{What it does not imply.} This is the region most exposed to
overreading, so the negative space is listed at length. \texttt{ACTIVE} does
not imply \texttt{INTENDED}: the magnitude of the projection can track the
dose perfectly while the induced activation moves off the registered axis or
against it (Section~\ref{ch07:sec:independence}). \texttt{ACTIVE} does not
imply causal efficacy of the driven content on behavior; it is a statement
about a projection coordinate, not about outputs. \texttt{ACTIVE} does not
imply that the linearity extends beyond the bracketing points, that the
region would survive a different seed or corpus, or that the fitted
resistance transfers to any other test point, model revision, or profile.
And \texttt{ACTIVE} at one interior index does not imply \texttt{ACTIVE}
anywhere else on the same sweep: region labels are pointwise.

\subsection{\texttt{SATURATION}}\label{ch07:sec:region-saturation}

\emph{Signature.} The backward slope is alive (above the slope floor) while
the forward slope is flat (at or below it): response grew on the way in and
has stopped growing on the way out. The implemented test is deliberately
one-sided and local; it detects a ceiling entered between the predecessor and
successor, not the global supremum of the response.

\emph{Causes.} Genuine ceilings in the projected activation; normalization
ceilings (a full-scale chosen low, so that the instrument clips before the
system does); representational ceilings such as norm growth limits at the
hooked site; or a sweep whose upper drives simply left the calibrated alpha
range's informative interior. The label does not distinguish a system that
saturates from an instrument that clips, and the candidate's separation of
the overload ratio from the clipped display value
(Chapter~\ref{ch08:quantization}) exists exactly so that this ambiguity is
carried in the record rather than erased by it.

\emph{What it does not imply.} \texttt{SATURATION} does not imply
\emph{oversaturation} in any semantic or safety sense; does not imply that
the response at the ceiling is aligned with the intended direction; does not
imply that dose reduction would restore linear response (hysteresis and state
dependence are explicitly out of scope of the memoryless candidate); and does
not imply that the ceiling is a property of the model rather than of the
measurement chain. Saturated-stage residuals are, per the source audit
itself, meaningless; the no-diagnosis machinery enforces this by refusing
residual computation outside \texttt{ACTIVE}.

\subsection{\texttt{BREAKDOWN}}\label{ch07:sec:region-breakdown}

\emph{Signature.} None is implemented numerically. \texttt{BREAKDOWN} is
retained in the enumeration from the source material --- the Guide's fifth
region, in which response ceases to be governed by the modeled relation at
all, by analogy with dielectric breakdown --- but the three-point classifier
has no test that emits it, and the audit's five-region summary and the
implementation's seven-member enumeration are both recorded in the seed paper
for exactness. As implemented, a \texttt{BREAKDOWN} label can enter the
system only as an \texttt{EXTERNAL\_UNVERIFIED} annotation.

\emph{Causes.} Hypothetically: interventions large enough to destroy the
local structure being measured --- degenerate outputs, representation
collapse, drives far beyond the calibrated grid. The candidate takes no
position on when this occurs, because it has no admitted evidence that it
occurs, and the calibrated alpha range gate refuses predictions outside the
grid in any case.

\emph{What it does not imply.} An external \texttt{BREAKDOWN} label carries
no numeric evidence by construction (the external mode forbids embedded
sweeps), so it implies nothing quantitative whatsoever. It is a typed way of
recording that some outside process asserted catastrophic non-response,
pending verification that the candidate itself cannot supply. It does not
imply damage, harm, jailbreak, or any behavioral state; the electrical
connotation of the borrowed word licenses nothing.

\subsection{\texttt{NON\_OHMIC}}\label{ch07:sec:region-nonohmic}

\emph{Signature.} Flow is resolvable and the point is not saturated, but the
relative slope error against the expected resistance exceeds the relative
tolerance: the response is doing something, and that something is not the
calibrated line. Convex, concave, kinked, and merely miscalibrated local
shapes all land here.

\emph{Causes.} Genuine nonlinearity of the projection in the swept range; a
calibrated resistance that is simply wrong for this neighborhood (fit
elsewhere, drifted, or fit on contaminated data); interference between the
steering intervention and other active mechanisms; or tolerance
misconfiguration. The label is a residual category in the taxonomic sense:
it collects everything resolvable that neither matches the line nor exhibits
the specific saturation signature.

\emph{What it does not imply.} \texttt{NON\_OHMIC} does not imply
\texttt{MISALIGNED}; this is one of the two implications the seed paper
singles out for explicit denial. A response can be strongly nonlinear in
magnitude while moving exactly along the intended axis, and the mnemonic
overlay accordingly returns no label at all for an intended non-ohmic point.
It also does not imply that the Ohmic model has been falsified in any global
sense --- one point exceeding one tolerance under one profile is not a test
of the registered hypotheses, whose closure protocol is preregistered and has
not been run.

\subsection{\texttt{UNSTABLE}}\label{ch07:sec:region-unstable}

\emph{Signature.} None is implemented numerically; like \texttt{BREAKDOWN},
\texttt{UNSTABLE} is reachable only as an external annotation. Its intended
meaning, inherited from the sources, is non-reproducibility: the same dose at
the same point under the same profile fails to yield a stable response
distribution --- variance rather than shape is the pathology.

\emph{Causes.} Seed sensitivity, sampling noise dominating the projection,
context-dependent activation at the hooked site, or drift within a session.
The three-point classifier, which sees a single sweep, is structurally
incapable of assessing reproducibility, which is why it never emits this
label; assessing it would require repeated sweeps under controlled
randomization, which is exactly item eight of the closure experiment list
(bootstrap uncertainty and cross-split reproducibility), an experiment that
has not been admitted.

\emph{What it does not imply.} \texttt{UNSTABLE} does not imply that the
underlying system is stochastic in any interesting sense, that the direction
is wrong, or that the profile is invalid; it implies only that a stable
pointwise response summary was not obtainable by whatever external process
asserted the label. It is also a no-diagnosis condition: unstable points
emit nothing downstream.

\section{A worked derivation, and what an assessment binds}\label{ch07:sec:worked}

The region taxonomy is easiest to hold onto with the classifier's
arithmetic run once in the open. All numbers below are illustrative
constructions over a synthetic sweep; nothing is a measurement.

\begin{example}[Illustrative three-point classification]\label{ch07:ex:sweep}
Let the profile declare expected resistance $r_{\mathrm{eff}} = 2.0$
alpha-per-activation, absolute tolerance $\texttt{delta\_abs} = 0.05$,
relative tolerance $\texttt{delta\_relative} = 0.2$, off-target maximum
$0.3$, and residual floor $0.01$. Consider the ordered sweep
observations (drive, flow, off-axis): $(1.0,\, 0.50,\, 0.05)$,
$(2.0,\, 1.00,\, 0.08)$, $(3.0,\, 1.48,\, 0.10)$, classified at the
interior index. The expected slope is $1/2.0 = 0.5$. The backward slope
is $(1.00 - 0.50)/(2.0 - 1.0) = 0.50$; the forward slope is
$(1.48 - 1.00)/(3.0 - 2.0) = 0.48$; the relative slope error is
$\max(|0.50 - 0.5|, |0.48 - 0.5|)/0.5 = 0.04$, well inside $0.2$. The
current flow $1.00$ exceeds \texttt{delta\_abs}, so not \texttt{CUTOFF};
the forward slope $0.48$ exceeds the slope floor
$0.5 \times 0.2 = 0.1$, so not \texttt{SATURATION}; the relative slope
error is inside tolerance, so not \texttt{NON\_OHMIC}: region
\texttt{ACTIVE}. The semantic ratio is $0.08 / \max(1.00, 0.01) =
0.08 \le 0.3$: outcome \texttt{INTENDED}. Overlay: \texttt{ALIGNED}.

Now perturb the same sweep three ways. (a) Flows
$(0.98,\, 1.00,\, 1.01)$: the backward slope collapses to $0.02$ and
the forward to $0.01$, both under the slope floor --- but the
saturation test requires a \emph{live} backward slope with a flat
forward slope, and here both are flat while the flow is large, so the
cascade falls through to the relative-slope-error test, which reads
$\max(|0.02-0.5|,|0.01-0.5|)/0.5 = 0.98 > 0.2$: region
\texttt{NON\_OHMIC} --- a uniformly flat plateau at high flow, entered
before the observed window, is locally indistinguishable from a wrong
slope, and the classifier says the modest thing rather than the
narrative thing. (b) Flows $(0.50,\, 1.00,\, 1.02)$: backward slope
$0.50$ (live), forward slope $0.02 \le 0.1$ (flat): region
\texttt{SATURATION} --- the ceiling was entered \emph{inside} the
window, which is what the one-sided signature detects. (c) Off-axis
column $(0.30,\, 0.40,\, 0.45)$ with the original flows: the semantic
ratio becomes $0.40$, exceeding $0.3$: outcome \texttt{OFF\_AXIS},
region still \texttt{ACTIVE}, overlay \texttt{MISALIGNED} --- the
independence principle in one line, since nothing about the magnitude
cascade changed.
\end{example}

The example also shows the sense in which tolerances are
\emph{classification parameters}, not implementation detail. Halving
\texttt{delta\_relative} to $0.1$ reclassifies perturbation (a) no
differently but tightens the \texttt{ACTIVE} verdict of the base case
($0.04$ still passes); raising it to $1.0$ would absorb perturbation
(a) into \texttt{ACTIVE}. Every region label is therefore implicitly
subscripted by the tolerance profile that produced it, and the
candidate makes the subscript explicit: the tolerance block is part of
the calibration profile, the profile digest is bound into every
assessment, and two assessments with different tolerance digests are
formally incomparable even over identical observations.

What the assessment object binds, beyond the two axis labels, completes
the account. A numeric assessment is factory-only (direct construction
is rejected by token) and carries: the classification pair; the mode
(\texttt{NUMERIC\_DERIVED}); the profile's response-region method
identity; the calibration-profile digest; the registered test-point
digest; an evidence reference whose digest must equal the canonical
digest of the embedded observation tuple (at least three points, with
an interior index); the index and expected resistance; the profile's
control bindings; and a full operation binding (Chapter~\ref{ch09:networks}'s
Definition on aggregate binding enumerates the binding fields). At
consumption time the assessment is not trusted but \emph{replayed}: the
validator re-resolves the evidence artifact, requires byte equality
with the embedded observations, re-derives the classification from the
observations, index, expected resistance, and current profile
tolerances, and requires the re-derived pair to equal the stored pair.
A classification that cannot be reproduced from its own bound evidence
does not validate --- which is the concrete meaning of the earlier
statement that region labels fail loudly.

By contrast, an \texttt{EXTERNAL\_UNVERIFIED} assessment must carry
\emph{no} observations, no index, and no expected resistance --- the
type rejects an external label that smuggles derived-looking numbers
--- and consumers are required to treat it as an unverified annotation:
the diagnostic layer, in particular, replays numeric assessments and
refuses external ones outright. The two modes thus split exactly along
the loud/silent failure line: numeric assessments carry their own
disproof machinery; external ones carry a label and an explicit
admission that no machinery stands behind it.

\section{The semantic outcome axis as implemented}\label{ch07:sec:outcome}

\begin{definition}[Semantic outcome enumeration]\label{ch07:def:outcome}
The implemented enumeration \texttt{SemanticOutcome} has exactly four
members: \texttt{UNASSESSED}, \texttt{INTENDED}, \texttt{OPPOSITE},
\texttt{OFF\_AXIS}.
\end{definition}

The numeric deriver assesses the semantic axis from a single ratio: the
off-axis activation magnitude at the current point divided by the flow
magnitude (floored by the profile's residual floor to avoid division by
values compatible with zero). If that ratio exceeds the profile's off-target
maximum, the outcome is \texttt{OFF\_AXIS}; otherwise \texttt{INTENDED}. Two
consequences follow and are recorded candidly. First, the numeric path never
emits \texttt{OPPOSITE}: distinguishing anti-directional movement from
off-axis movement requires signed projection evidence that the magnitude-only
sweep observation does not carry, so \texttt{OPPOSITE} is, in the current
implementation, an externally assessable outcome only. Second, the numeric
path never emits \texttt{UNASSESSED}, because it refuses to run at all
without the observations that make the ratio computable; \texttt{UNASSESSED}
is what an assessment carries when the semantic question was not put to any
evidence.

\subsection{\texttt{UNASSESSED}}

The typed absence of a directional judgment. A profile can classify region
crisply and still have no admissible answer to ``did it move the intended
way,'' most obviously when no behavioral or activation-space comparison
against the registered direction has been performed. \texttt{UNASSESSED} on
the semantic axis is a hard silence condition: the residual diagnostic
explicitly rejects assessments whose semantic outcome is \texttt{UNASSESSED},
on the principle that a residual against an intention nobody assessed is an
invitation to invent the intention afterward.

\subsection{\texttt{INTENDED}}

The observed activation moved along the registered direction, with off-axis
contamination below the profile's tolerance. \texttt{INTENDED} does not
imply that the magnitude was appropriate (it combines with every region,
including \texttt{CUTOFF} and \texttt{SATURATION}), does not imply that the
behavioral consequence of the activation movement matches anyone's intention
at the output level, and does not imply endorsement of the intervention. The
word ``intended'' binds to the registered direction vector and to nothing
else: change the registration and the same physical event reads differently.
This is deliberate; intention lives in the profile, where it is frozen,
digested, and auditable, not in the interpreter's retrospect.

\subsection{\texttt{OPPOSITE}}

The observed activation moved against the registered direction: a signed
reversal, not merely leakage. As noted, the current numeric path cannot
produce this label; it exists in the enumeration because the source material
treats reversal as a first-class outcome (the inverted-response reading of a
steering intervention) and because a future signed-evidence classifier must
have somewhere typed to land. An \texttt{OPPOSITE} label, where externally
asserted, does not imply psychological reactance, adversarial optimization,
safety-training interference, or any mechanism at all; mechanism attribution
from outcome labels is exactly the inferential step the independence
principle below forbids.

\subsection{\texttt{OFF\_AXIS}}

The observed activation moved, but a supra-tolerance fraction of the
movement lies outside the registered direction. This is the label the
numeric path emits when the off-axis ratio exceeds the profile's
off-target maximum. \texttt{OFF\_AXIS} does not imply that the off-axis
component is harmful, meaningful, or even stable; a direction vector that
captures its target imperfectly will manufacture off-axis readings from
perfectly on-target mechanisms, and a poorly conditioned projection basis
will do the same. The label records a geometric fact about one comparison
under one registration; the temptation to read it as ``side effects'' is a
semantic overlay the candidate declines.

\section{The independence principle}\label{ch07:sec:independence}

\begin{ruleenv}[Axis independence]\label{ch07:rule:independence}
No dose-response region mechanically determines a semantic outcome, and no
semantic outcome mechanically determines a region. Every pair in the product
$\{$region$\} \times \{$outcome$\}$ is admissible to the type system, and the
classification container stores the two axes as separate fields with no
consistency coupling beyond enumeration membership.
\end{ruleenv}

The principle is worth developing at length, because it is the load-bearing
wall between this chapter and the mnemonic it replaces.

\emph{Active does not mean aligned.} The region axis sees the magnitude of a
projection; the outcome axis sees its geometry. A dose that recruits a
strong, beautifully linear response in the projection coordinate can be
recruiting it off the intended axis --- the projection tracks whatever
component of the induced activation lies along the measurement direction,
and a diagonal movement has such a component. Conversely, an intervention
can move activation precisely along the intended axis while the magnitude
response is grossly nonlinear. In the crossed cases, the mnemonic has no
word, and that lexical gap in ASOM is not an accident of vocabulary: it
marks a real, observable class of events that a four-word overlay cannot
name. An \texttt{ACTIVE}/\texttt{OFF\_AXIS} classification and an
\texttt{ACTIVE}/\texttt{OPPOSITE} classification are the two canonical
counterexamples to region-implies-outcome, and both are fully admissible
records.

\emph{Flat response is multi-causal.} The converse direction is subtler and
more dangerous, because flatness invites narrative. A flat or negligible
response at the classified point is evidence-compatible with at least five
generating situations, which the candidate insists on keeping distinct:
(i) \emph{cutoff} --- the dose is below whatever threshold governs
resolvable response at this point; (ii) \emph{saturation} --- the response
exhausted its range before the observed doses, so the local curve is flat at
ceiling rather than at floor; (iii) \emph{wrong direction} --- the
intervention moves activation vigorously, but orthogonally to the
measurement direction, so the projection sees nothing; (iv)
\emph{insensitivity} --- the system genuinely does not respond to this
intervention at this site, at any dose; and (v) \emph{measurement failure}
--- hooks misplaced, units mismatched, full-scale misdeclared, instrument
silent. Of these five, the three-point classifier can discriminate only the
first two, and only by their local slope signatures; the last three all
present as \texttt{CUTOFF} or as low-magnitude \texttt{ACTIVE} noise, and
distinguishing them requires controls --- random same-norm directions,
wrong-layer placements, zero-intervention baselines --- that live in the
preregistered closure protocol, not in the classifier. A flat point is
therefore never, by itself, evidence of starvation, and the candidate's
refusal to emit \textsc{starved} except through the explicitly decorative
overlay is the operational form of that refusal.

The independence principle also has an evidential asymmetry worth recording.
Region classification is internally checkable: it replays from the
observation tuple and the tolerance profile, and the factory binds the digest
of the observations so that replay is exact. Outcome classification, by
contrast, is only as good as the direction registration and the off-axis
instrumentation, both of which are profile inputs that the classifier
consumes on trust. The two axes therefore fail differently --- region labels
fail loudly, by replay divergence, while outcome labels fail silently, by
mis-registration --- and treating them as one axis would let the silent
failure mode contaminate the loud one. Keeping them separate keeps their
error budgets separate, in the same spirit in which model verification and
model validation are kept separate activities \cite{sargent}.

\section{ASOM as mnemonic overlay}\label{ch07:sec:asom}

\begin{definition}[ASOM overlay]\label{ch07:def:asom}
The implemented enumeration \texttt{AsomLabel} has four members:
\texttt{STARVED}, \texttt{ALIGNED}, \texttt{OVERSATURATED},
\texttt{MISALIGNED}. The overlay is the partial function
\texttt{asom\_label()} on classifications, defined exactly as follows:
if the semantic outcome is \texttt{OPPOSITE} or \texttt{OFF\_AXIS}, return
\texttt{MISALIGNED}; otherwise, if the semantic outcome is not
\texttt{INTENDED}, return nothing; otherwise return \texttt{STARVED} for
region \texttt{CUTOFF}, \texttt{ALIGNED} for \texttt{ACTIVE},
\texttt{OVERSATURATED} for \texttt{SATURATION}, and nothing for every other
region.
\end{definition}

Three properties of this definition carry the argument. First, the overlay
is \emph{derived}: it is computed from the two stored axes and stored
nowhere, so no ASOM word can drift out of agreement with the axes that
justify it, and no process can write an ASOM word directly. Second, the
overlay is \emph{partial}: an \texttt{UNASSESSED} outcome yields no label at
all, and an \texttt{INTENDED} outcome over \texttt{BREAKDOWN},
\texttt{NON\_OHMIC}, \texttt{UNSTABLE}, or \texttt{UNASSESSED} regions
likewise yields nothing --- the mnemonic simply has no word for most of the
product space, and the implementation returns \texttt{None} rather than
coining one. Third, the overlay is \emph{lossy in exactly one direction}:
\texttt{MISALIGNED} merges \texttt{OPPOSITE} with \texttt{OFF\_AXIS} and
discards the region entirely, so recovering the axes from the label is
impossible. A four-word summary of a twenty-eight-cell product space cannot
be information-preserving, and the candidate's position is that it should
not try: the words exist for human pattern recognition in displays and
prose, and every load-bearing consumer reads the axes.

\begin{ruleenv}[No identity theorem]\label{ch07:rule:asom-no-identity}
No inference, gate, capability, diagnostic, or claim in the candidate
consumes an \texttt{AsomLabel}. The overlay never appears in canonical
payloads of measurement values, never conditions residual admission, and
never participates in packet completeness. Any external document that
reasons from \textsc{starved}, \textsc{aligned}, \textsc{oversaturated}, or
\textsc{misaligned} as if these were classification primitives is reasoning
from a display artifact.
\end{ruleenv}

\section{The full region-by-outcome cross table}\label{ch07:sec:crosstable}

Table~\ref{ch07:tab:cross} enumerates the product space. For compactness
the outcome columns are abbreviated U (\texttt{UNASSESSED}),
I (\texttt{INTENDED}), Op (\texttt{OPPOSITE}), Off (\texttt{OFF\_AXIS});
each cell records the derived ASOM overlay, if any, with ``---'' for the
overlay's silence.

\begin{table}[htbp]
\centering
\begin{tabular}{@{}lcccc@{}}
\toprule
Region & U & I & Op & Off \\
\midrule
\texttt{UNASSESSED} & --- & --- & \texttt{MISALIGNED} & \texttt{MISALIGNED} \\
\texttt{CUTOFF} & --- & \texttt{STARVED} & \texttt{MISALIGNED} & \texttt{MISALIGNED} \\
\texttt{ACTIVE} & --- & \texttt{ALIGNED} & \texttt{MISALIGNED} & \texttt{MISALIGNED} \\
\texttt{SATURATION} & --- & \texttt{OVERSATURATED} & \texttt{MISALIGNED} & \texttt{MISALIGNED} \\
\texttt{BREAKDOWN} & --- & --- & \texttt{MISALIGNED} & \texttt{MISALIGNED} \\
\texttt{NON\_OHMIC} & --- & --- & \texttt{MISALIGNED} & \texttt{MISALIGNED} \\
\texttt{UNSTABLE} & --- & --- & \texttt{MISALIGNED} & \texttt{MISALIGNED} \\
\bottomrule
\end{tabular}
\caption{The $7 \times 4$ classification product and its derived ASOM
overlay. The overlay names only six of the twenty-eight cells distinctly;
fourteen cells share \texttt{MISALIGNED}, and eight cells have no name.}
\label{ch07:tab:cross}
\end{table}

The remainder of this section walks the table row by row, reading each
combination for what it can and cannot mean. Throughout, ``possible
readings'' are enumerations of evidence-compatible situations, not
diagnoses; selecting among them is exactly what classification alone cannot
do.

\emph{\texttt{UNASSESSED} region.} With outcome U, this is the null
assessment: a typed placeholder, common for freshly registered test points.
With outcome I, an external process asserted directional success without any
region evidence --- admissible as a record, useless as a premise; the
overlay pointedly has no word for it. With Op or Off, someone asserted
directional failure with no dose-response context; the overlay says
\texttt{MISALIGNED}, and the correct reading of that word here is ``an
unverified external annotation exists,'' nothing stronger.

\emph{\texttt{CUTOFF} region.} With U: no resolvable flow and no
directional assessment --- the modal honest state for a dead point, since
assessing direction of a sub-tolerance response is numerically fraught.
With I: the overlay's \texttt{STARVED} cell; the response is below
resolution and whatever trace exists is on-axis. Possible readings include
genuine sub-threshold dosing, but equally a dead instrument whose noise
floor happens to project on-axis; the cell does not separate them. With Op:
below-resolution flow asserted to be reversed --- a nearly self-undermining
combination that should prompt scrutiny of the asserting process rather than
of the model. With Off: sub-tolerance magnitude, off-axis geometry; often
the signature of measuring the wrong thing in the right place, and a natural
trigger for the wrong-direction hypothesis of
Section~\ref{ch07:sec:independence}, which only controls can confirm.

\emph{\texttt{ACTIVE} region.} With U: a calibrated-looking response whose
direction nobody has assessed; the diagnostic machinery refuses to run on
it, precisely because an unassessed intention invites retrofitting. With I:
the overlay's \texttt{ALIGNED} cell and the candidate's best case ---
magnitude tracks dose at the calibrated slope, geometry tracks the
registration; still a structural statement under one profile, not an
endorsement, not transferable, not causal. With Op: linear-in-magnitude,
reversed-in-direction; if externally asserted and later verified, this
combination would be the single most alarming cell for a steering
application, and the candidate's contribution is to make it representable
rather than lexically impossible (ASOM has no word for a linear reversal
other than the undifferentiated \texttt{MISALIGNED}). With Off: strong,
lawful magnitude response leaking off-axis --- the canonical
``active-but-off-axis'' counterexample; possible readings span imperfect
direction registration, genuinely diffuse mechanisms, and superposed
responses, and the cell selects none of them.

\emph{\texttt{SATURATION} region.} With U: ceiling behavior, direction
unassessed; residuals refuse to run, and per the sources saturated-stage
residuals would be meaningless anyway. With I: the overlay's
\texttt{OVERSATURATED} cell --- response climbed on-axis and hit a ceiling;
readings include a genuine representational ceiling and a clipping
instrument, which the overload ratio (kept separately from the clipped
display, Chapter~\ref{ch08:quantization}) exists to help distinguish. With
Op: ceiling-flat magnitude with asserted reversal --- coherent only if the
reversal occurred below the ceiling or in a signed coordinate the sweep does
not carry; treat as external assertion pending signed evidence. With Off:
the ceiling is real but the content at ceiling is off-target; a reading
consistent with saturating the wrong mechanism entirely.

\emph{\texttt{BREAKDOWN} region.} All four cells are external-only, since
the numeric path never emits the region. With U, the cell is a bare
catastrophic-response assertion; with I, an internally tense one (broken
response, yet on-axis?); with Op or Off, the overlay's undifferentiated
\texttt{MISALIGNED} again. In every case the assessment carries no
observations by construction, and the only defensible consumption of the
row is as a flag for human review of the asserting process.

\emph{\texttt{NON\_OHMIC} region.} With U: resolvable, non-linear
magnitude response, direction unassessed --- an honest and common state for
exploratory sweeps. With I: nonlinear but on-axis; the overlay has no word,
and rightly so --- ``misaligned'' would be false and ``aligned'' would
launder the nonlinearity. Readings include genuine supra-linear
recruitment, miscalibrated expected resistance, and interference; the
falsification program's gain-composition and inverse-slope experiments are
the registered instruments for narrowing this set. With Op or Off: nonlinear
magnitude and adverse geometry; the least constrained cells in the table,
compatible with nearly any mechanism story and therefore probative of none.

\emph{\texttt{UNSTABLE} region.} All cells are external-only. With U: an
assertion of irreproducibility, unquantified. With I: ``unstable but
on-average intended'' --- a reading that presupposes distributional evidence
the record cannot carry, hence flagged, not consumed. With Op or Off: the
overlay says \texttt{MISALIGNED}, but the honest reading is that neither
axis is trustworthy at an unstable point, which is why instability is a
no-diagnosis condition rather than a diagnosable state.

\begin{remark}\label{ch07:rem:table-use}
The table's practical function is negative. It exists so that no
combination can be declared impossible by vocabulary, and so that every
overlay word can be traced to the cells it covers and the cells it merges.
Twenty-two of twenty-eight cells receive either no ASOM word or a shared
one; any analysis pipeline that operates at ASOM granularity therefore
discards most of the classification structure by construction.
\end{remark}

\section{No-diagnosis conditions}\label{ch07:sec:nodiag}

The candidate's diagnostic layer (residual computation, and everything
downstream of it) is guarded by a conjunction of admission conditions, each
of which independently silences diagnosis. They are enumerated here as
implemented, with the enforcing mechanism named for each.

\begin{ruleenv}[Silence conditions]\label{ch07:rule:nodiag}
No diagnostic is emitted when any of the following holds:
\begin{enumerate}
\item \emph{Region is not \texttt{ACTIVE}.} The residual entry point
rejects any assessment whose region differs from \texttt{ACTIVE}; this
covers \texttt{CUTOFF}, \texttt{SATURATION}, \texttt{BREAKDOWN},
\texttt{NON\_OHMIC}, \texttt{UNSTABLE}, and \texttt{UNASSESSED} uniformly.
Cutoff points have nothing to diagnose; saturated points have meaningless
residuals; unstable points have untrustworthy ones.
\item \emph{Semantic outcome is \texttt{UNASSESSED}.} A residual against an
unassessed intention is refused outright.
\item \emph{Profile is uncalibrated or not reviewed.} The operation context
admits a calibration profile only in review state \texttt{REVIEWED}; any
other review state fails the context, so no operation --- classification
included --- can even begin.
\item \emph{Profile drift state is not current.} A profile whose drift
state differs from current is rejected at context admission.
\item \emph{Replay time is stale.} The caller-supplied \texttt{replay\_as\_of}
must satisfy $\texttt{valid\_from} \le \texttt{replay\_as\_of} <
\texttt{expires\_at}$; outside that window the profile is ``stale or not yet
valid'' and the context fails.
\item \emph{Wrong model or wrong profile.} Every assessment binds the
calibration profile digest and a registered test-point digest; an
assessment presented against a context whose profile digest differs, or
whose test point is unregistered, is rejected. Since the profile binds the
exact model, configuration, tokenizer, and weights identities, a wrong-model
presentation is a wrong-profile presentation and fails the same check.
\end{enumerate}
\end{ruleenv}

Two design commitments underlie the list. The first is that silence is the
default and diagnosis is the exception: every condition is checked
fail-closed, and a missing precondition produces a typed rejection, never a
degraded diagnostic. The second is that the conditions compose by
conjunction with no override path --- there is no force flag, no
administrative bypass, and no trust mode in which a stale profile or an
unassessed outcome becomes diagnosable. The alternative --- emitting
diagnostics with warnings attached --- was foreclosed by the source audit's
observation that warned-against readings (the residual-as-bypass reading in
particular) were made on the same page as the warnings.

\section{Replay-time semantics: \texttt{replay\_as\_of}}\label{ch07:sec:replay}

Every operation context carries a field \texttt{replay\_as\_of}: a
caller-supplied timestamp, parsed and canonicalized at admission, snapshotted,
and immutable thereafter (any change after admission is a validation error).
Profile currency (Rule~\ref{ch07:rule:nodiag}, condition 5) is evaluated
against this timestamp, and every operation binding embeds it, so that each
result records the historical instant \emph{as of which} the profile was
asserted current.

\begin{ruleenv}[Replay time is not authenticated freshness]\label{ch07:rule:replay}
\texttt{replay\_as\_of} is a replay coordinate, not a clock. It asserts that
the computation is being evaluated as of a stated historical time; it does
not assert that the stated time is now, that the caller's clock is honest,
or that the profile is current in wall-clock reality. The context exposes
this directly: its \texttt{authenticated\_current} property is constantly
false, and the optional clock-receipt field, when present, is itself only a
recorded artifact, not a trusted time source. A caller can therefore always
choose a \texttt{replay\_as\_of} inside a profile's validity window and
obtain classifications against that historical frame; nothing in the
candidate prevents this, and nothing in the candidate represents the result
as fresh.
\end{ruleenv}

The design is deliberate and mirrors the completeness-versus-eligibility
split that governs the whole candidate: structural checks (the window
inequality, the snapshot immutability, the binding of the timestamp into
every result) are implemented and tested; authority checks (a trusted
current-time receipt from an admitted external authority) are named,
required for authenticated Evidence-Packet eligibility, and absent. The
honest reading of any classification record is therefore always
conditional: \emph{under profile $P$, replayed as of $t$, this sweep
classified as such}. Removing the conditionality would require exactly the
authenticated infrastructure the candidate declares it does not have.

\begin{remark}[What this chapter establishes]\label{ch07:rem:close}
The classification layer contributes no empirical claim. What it
contributes is a shape for future claims: a seven-member region axis and a
four-member outcome axis, independent by rule, with a decorative overlay
that cannot be consumed, a product space in which no combination is
unspeakable, and a silence discipline under which most states of the world
produce no diagnosis at all. If the registered hypotheses eventually close
against this scaffolding, their content will be expressible without
renaming anything; if they fail, the failure will be recordable in the same
vocabulary. That neutrality between success and failure is the property the
chapter was built to secure.
\end{remark}
